\begin{solapartat}
\punts{0,4} El Sol està en el focus de l'el·lipse i a l'hivern la Terra s'ha de trobar al periheli o a prop d'aquest.
\figura[0.5]{sol_fig1.png}

\punts{0,4} L'energia mecànica de la Terra es conserva atès que només actua la força de la gravetat, que és una força conservativa.
L'energia mecànica és la suma de les energies cinètica i potencial gravitatòria:
\[ E_m(r) = \frac{1}{2}M_T v^2 - G\,\frac{M_T M_S}{r} = \mathit{constant} \]
Quan $r$ és mínim, llavors el terme $G\,\frac{M_S M_T}{r}$ és màxim i, per tant, l'energia potencial és mínima al periheli i màxima a l'afeli.
Com que l'energia mecànica és constant, si a l'afeli l'energia potencial gravitatòria és màxima, llavors l'energia cinètica serà mínima i també serà mínima la velocitat en aquest punt de la trajectòria.

També es pot argumentar que:
\[ E_m(r) = \frac{1}{2}M_T v^2 - G\,\frac{M_T M_S}{r} \;\Longrightarrow\; v = \sqrt{\mathit{constant} + 2G\,\frac{M_S}{r}} \]
Com que a l'afeli $r$ és màxim, llavors el terme $2G\,\frac{M_S}{r}$ és mínim i també ho serà la velocitat.

Es pot plantejar la solució a partir de la segona llei de Kepler:

\punts{0,35} $r_a v_a M_T = r_p v_p M_T$

\punts{0,1} $r_a = r_p\,\dfrac{v_p}{v_a} = 147,1\times 10^{9}\,\dfrac{30.750}{28.760} = 1,57\times 10^{11}\ \mathrm{m}$

\destaca{Alternativament,} es pot plantejar a partir de la conservació de l'energia mecànica:

\punts{0,1} $E_m(r_a) = E_m(r_p)$

\punts{0,25} $\dfrac{1}{2}M_T v_a^2 - G\,\dfrac{M_T M_S}{r_a} = \dfrac{1}{2}M_T v_p^2 - G\,\dfrac{M_T M_S}{r_p}$
\[ \frac{1}{r_a} = \frac{1}{r_p} + \frac{1}{2}\,\frac{v_a^2 - v_p^2}{G M_S} \]
\punts{0,1} $r_a = \left[\dfrac{1}{r_p} + \dfrac{1}{2}\,\dfrac{v_a^2 - v_p^2}{G M_S}\right]^{-1}$
\[ r_a = \left[\frac{1}{147,1\times 10^{9}} + \frac{1}{2}\,\frac{28.760^2 - 30.750^2}{6,67\times 10^{-11}\cdot 1,99\times 10^{30}}\right]^{-1} = 1,57\times 10^{11}\ \mathrm{m} \]
\end{solapartat}

\begin{solapartat}
\punts{0,65} $g_S = G\,\dfrac{M_S}{R_S^2} = 274\ \mathrm{m/s^2}$ \qquad \punts{0,1}

\punts{0,4} El pes és: \quad $\mathit{Pes} = m\,g_S = 274\cdot 10,0 = 2740\ \mathrm{N}$ \qquad \punts{0,1}
\end{solapartat}
